% Experimental-design review v1 (Portuguese report).
% Build instructions are in README.md of this directory.
\documentclass[11pt,a4paper]{article}

\usepackage[T1]{fontenc}
\usepackage{lmodern}
\usepackage{babel}
% The local TeX installation has no Portuguese hyphenation patterns, so the
% text is set without hyphenation instead of with foreign patterns.
\babelprovide[import,main,hyphenrules=nohyphenation]{portuguese}
\usepackage[a4paper,margin=2.3cm,headheight=14pt]{geometry}
\usepackage{microtype}
\usepackage{amsmath,amssymb}
\usepackage{booktabs,tabularx,array,multirow,longtable}
\usepackage[table]{xcolor}
\usepackage{tikz}
\usepackage{standalone}
\usepackage{graphicx}
\usepackage{pdflscape}
\usepackage[font=small,labelfont=bf,labelsep=period]{caption}
\usepackage{enumitem}
\usepackage[most]{tcolorbox}
\usepackage{ragged2e}
\usepackage{fancyhdr}
\usepackage[round,authoryear]{natbib}
\usepackage[hidelinks]{hyperref}

% Shared TikZ style for the report figures.
% Colours follow the Okabe-Ito palette, which stays distinguishable under the
% common forms of colour-vision deficiency. Every colour-coded element also
% carries a text label, so no figure depends on colour alone.
\usetikzlibrary{arrows.meta,positioning,fit,calc,backgrounds,decorations.pathreplacing,shapes.geometric}

\definecolor{oiOrange}{HTML}{E69F00}
\definecolor{oiSky}{HTML}{56B4E9}
\definecolor{oiGreen}{HTML}{009E73}
\definecolor{oiYellow}{HTML}{F0E442}
\definecolor{oiBlue}{HTML}{0072B2}
\definecolor{oiVerm}{HTML}{D55E00}
\definecolor{oiPurple}{HTML}{CC79A7}
\definecolor{oiGray}{HTML}{6E6E6E}

\colorlet{corTreino}{oiBlue!22}
\colorlet{corValid}{oiGreen!30}
\colorlet{corTeste}{oiOrange!45}
\colorlet{corControle}{oiPurple!35}

\tikzset{
  >=Stealth,
  bloco/.style={draw=black!70, rounded corners=1.5pt, minimum height=0.62cm,
    inner sep=2pt, font=\sffamily\fsz, align=center},
  treino/.style={bloco, fill=corTreino},
  valid/.style={bloco, fill=corValid},
  teste/.style={bloco, fill=corTeste, thick},
  etapa/.style={draw=black!70, rounded corners=3pt, fill=black!4,
    text width=4.3cm, align=center, inner sep=4pt, font=\sffamily\fsz},
  nota/.style={font=\sffamily\fsz, align=left, text=black!80},
  titulo/.style={font=\sffamily\fst\bfseries, anchor=west},
  seta/.style={->, thick, black!70},
  regiao/.style={draw=black!50, dashed, rounded corners=4pt},
}

% Absolute font sizes, so a figure looks the same standalone (10 pt class) and
% inside the 11 pt report.
\providecommand{\fsx}{\fontsize{5.5}{6.5}\selectfont}
\providecommand{\fsz}{\fontsize{7}{8.5}\selectfont}
\providecommand{\fst}{\fontsize{9}{11}\selectfont}


\setlength{\emergencystretch}{3em}
\setlength{\parskip}{0.35em}
\setlength{\parindent}{0pt}
\setlist{itemsep=0.15em, topsep=0.3em}
\renewcommand{\arraystretch}{1.18}
\newcolumntype{Y}{>{\RaggedRight\arraybackslash}X}
\newcolumntype{L}[1]{>{\RaggedRight\arraybackslash}p{#1}}

\pagestyle{fancy}
\fancyhf{}
\fancyhead[L]{\small\sffamily CapsProm: avaliação do desenho experimental}
\fancyhead[R]{\small\sffamily v1, 4 nov.\ 2026}
\fancyfoot[C]{\small\thepage}

% Evidence tags. Each tag carries its own word, so colour is not needed to read it.
\newcommand{\etiqueta}[2]{\tcbox[on line, colback=#1!10, colframe=#1!75!black,
  boxrule=0.4pt, arc=1.2pt, left=1.5pt, right=1.5pt, top=0pt, bottom=0pt,
  fontupper=\sffamily\scriptsize\bfseries]{#2}}
\newcommand{\lit}{\etiqueta{oiBlue}{Literatura}}
\newcommand{\proj}{\etiqueta{oiGreen}{Projeto}}
\newcommand{\rec}{\etiqueta{oiVerm}{Recomendação}}
\newcommand{\fic}{\etiqueta{oiPurple}{Fictício}}
\newcommand{\arq}[1]{\nolinkurl{#1}}

\newtcolorbox{exemplo}[1]{breakable, enhanced, colback=oiPurple!4, colframe=oiPurple!70!black,
  boxrule=0.6pt, arc=2pt, title={Exemplo fictício: #1}, fonttitle=\sffamily\bfseries\small}
\newtcolorbox{resumo}[1][Em resumo]{breakable, enhanced, colback=oiSky!7, colframe=oiBlue!65!black,
  boxrule=0.6pt, arc=2pt, title={#1}, fonttitle=\sffamily\bfseries\small}

\newcommand{\MCC}{\mathrm{MCC}}
\DeclareMathOperator{\Var}{Var}
\DeclareMathOperator{\med}{mediana}

\title{\sffamily\bfseries Avaliação crítica de um desenho experimental aninhado para o CapsProm\\[4pt]
\large\mdseries $K$-fold repetido com teste $t$ corrigido no nível externo e 5$\times$2cv com o $F$ combinado de Alpaydin no nível interno}
\author{\sffamily Relatório metodológico para a tese de doutorado em Ciência da Computação (UFOP)\\
\sffamily\small Projeto CapsProm\quad·\quad elaborado com assistência do Claude (Anthropic), para revisão do autor}
\date{\sffamily\small Versão 1, 4 de novembro de 2026}

\begin{document}
\maketitle
\thispagestyle{fancy}

\begin{resumo}[Resumo executivo]
\textbf{Pergunta.} Convém adotar, no CapsProm, um desenho em dois níveis: $K$-fold repetido com $S$ sementes e teste $t$ corrigido no nível externo, e 5$\times$2cv com o teste $F$ combinado de Alpaydin no nível interno?

\textbf{Resposta curta.} Não como desenho confirmatório principal. Partes dele são úteis em papéis secundários.
\begin{enumerate}[leftmargin=1.4em]
\item \textbf{O nível interno não precisa de teste de hipótese.} Esse nível serve para \emph{escolher} uma configuração, e escolher é uma decisão, não uma inferência. Usar o $F$ de Alpaydin para escolher traz vários problemas:
  \begin{itemize}[leftmargin=1.2em]
  \item multiplica os testes;
  \item confunde ``não significativo'' com ``equivalente'';
  \item treina cada modelo com só metade do treino externo;
  \item custa o dobro de um 5-fold interno.
  \end{itemize}
  Uma regra de seleção fixada de antemão resolve o problema por menos.
\item \textbf{O nível externo com teste $t$ corrigido é defensável, com limites.} Ele estima o desempenho de um \emph{procedimento} de aprendizagem, não do modelo final ajustado. A correção de Nadeau e Bengio é heurística. O teste vale para um dataset por vez e exige folds formados por grupos de homologia. Ele não substitui a inferência entre datasets, que no projeto se faz por cluster de dependência.
\item \textbf{A proposta conflita com o protocolo vigente.} Um $K$-fold externo reciclaria o conjunto de teste externo selado, que hoje é fixo, separado por grupos e aberto uma única vez. Para o $K$-fold repetido não conflitar, ele deve ficar \emph{dentro} da partição de treino, como análise de estabilidade.
\item \textbf{Controles negativos pedem outro desenho.} O efeito da arquitetura, o efeito da estratégia de controles e a interação entre os dois exigem um desenho fatorial avaliado num painel de teste comum. Também exigem diagnósticos de qualidade dos controles: um MCC maior não indica controle melhor.
\end{enumerate}
\textbf{Recomendação.} Manter o protocolo confirmatório vigente (\S\ref{sec:recomendacao}). Trocar qualquer teste interno por uma regra de seleção predeclarada. Acrescentar como análises secundárias: (i) $K$-fold por grupos repetido dentro do treino, para estabilidade; (ii) 5$\times$2cv apenas em experimentos dedicados a um par de procedimentos; (iii) o desenho fatorial de controles, com diagnósticos.
\end{resumo}

\paragraph{Como ler as etiquetas.} Cada afirmação importante traz sua origem:
\begin{itemize}[leftmargin=1.4em]
\item \lit\ resultado publicado em fonte primária;
\item \proj\ fato verificado nos arquivos do repositório nesta data, com o caminho do arquivo;
\item \rec\ proposta deste relatório, que ainda não foi adotada;
\item \fic\ número inventado só para ilustrar, sem relação com nenhum resultado do projeto.
\end{itemize}
Este relatório não altera o protocolo vigente e não inicia experimentos.

\tableofcontents

%=====================================================================
\section{Escopo e fontes consultadas}
%=====================================================================

O relatório examina a proposta abaixo sem presumir que ela seja adequada:

\begin{itemize}[leftmargin=1.4em]
\item \textbf{Nível externo:} validação cruzada $K$-fold repetida com $S$ sementes, associada a um teste $t$ corrigido para reamostragem \citep{nadeau2003,bouckaert2004}.
\item \textbf{Nível interno:} 5$\times$2cv associado ao teste $F$ combinado de \citet{alpaydin1999}.
\end{itemize}

\proj\ Foram lidos os seguintes arquivos do repositório:
\begin{itemize}[leftmargin=1.4em]
\item o protocolo estatístico, \arq{specs/statistical_evaluation_protocol.md}, e seu contrato vigente, \arq{configs/evaluation_protocol_v4.yaml} (versão 4.0.0, desenho congelado e execução bloqueada);
\item a política de sementes, \arq{specs/seed_policy.md};
\item os protocolos de controles negativos (\arq{specs/negative_control_protocol.md} e \arq{specs/cds_negative_control_protocol.md}) e de agrupamento por homologia (\arq{specs/homology_grouping_protocol.md});
\item o contrato do experimento canônico, \arq{specs/capsprom_canonical_experiment.md};
\item o contrato de reconstrução genômica, \arq{configs/graphpro_reconstruction_v3.yaml};
\item o registro de achados, \arq{docs/provenance/findings_register.md}.
\end{itemize}

\lit\ As fontes estatísticas primárias são:
\begin{itemize}[leftmargin=1.4em]
\item \citet{dietterich1998}, \citet{alpaydin1999}, \citet{nadeau2003} e \citet{bouckaert2004};
\item \citet{bengio2004} e \citet{bates2024};
\item \citet{varma2006} e \citet{cawley2010};
\item \citet{demsar2006}, \citet{benavoli2017} e \citet{holm1979};
\item \citet{field2007}, \citet{teufel2023} e \citet{petti2022}.
\end{itemize}
A fórmula do teste corrigido para $K$-fold repetido foi conferida no texto completo de \citet{bouckaert2004}.

%=====================================================================
\section{Quatro atividades que não devem ser confundidas}
\label{sec:conceitos}
%=====================================================================

Um desenho experimental em aprendizado de máquina mistura quatro atividades. Cada uma responde a uma pergunta diferente e erra de um jeito diferente (Tabela~\ref{tab:conceitos}).

\begin{table}[htbp]
\centering\small
\caption{As quatro atividades de um desenho experimental. A confusão entre elas é a origem da maioria dos erros discutidos neste relatório.}
\label{tab:conceitos}
\begin{tabularx}{\linewidth}{L{2.6cm}YYY}
\toprule
Atividade & Pergunta respondida & Produto & Erro típico \\
\midrule
Reamostragem & Como gerar várias versões de treino e teste a partir de um dataset? & Partições (folds, metades, réplicas bootstrap) & Partir pares ou homólogos entre treino e teste (vazamento) \\
Seleção de modelos & Qual configuração usar? & Uma configuração $c^\ast$ (hiperparâmetros, épocas, limiar) & Escolher olhando o teste; tratar a escolha como teste de hipótese \\
Estimação de generalização & Quanto o procedimento acerta em dados novos da mesma população? & Estimativa pontual e intervalo & Estimar no mesmo dado usado para escolher (viés otimista) \\
Teste de hipóteses & A diferença observada entre A e B é compatível com o acaso? & Valor-$p$, decisão com erro tipo I controlado & Tratar folds ou sementes como observações independentes \\
\bottomrule
\end{tabularx}
\end{table}

\lit\ \citet{varma2006} e \citet{cawley2010} mostram que usar a mesma validação cruzada para escolher e para estimar produz estimativas otimistas. A validação cruzada aninhada separa as duas coisas: o laço interno escolhe e o laço externo estima.

\lit\ \citet{bates2024} mostram outra coisa. O que a validação cruzada estima é o erro médio de modelos treinados em \emph{outros} conjuntos de treino da mesma população, e não o erro do modelo ajustado nos dados em mãos. Além disso, intervalos de confiança ingênuos sobre essa estimativa cobrem menos do que o nível nominal.

%=====================================================================
\section{O desenho proposto}
\label{sec:desenho}
%=====================================================================

\subsection{Explicação intuitiva}

Imagine um concurso entre dois métodos de cozinha, A e B, julgado com os ingredientes disponíveis.

No \emph{nível externo}, separamos os ingredientes em $K$ lotes. Em cada rodada, um lote fica guardado para a prova final e os outros vão para a cozinha. Repetimos a separação $S$ vezes, sorteando os lotes de novo a cada vez.

No \emph{nível interno}, dentro da cozinha de cada rodada, cada método ainda precisa decidir sua receita: tempo de forno, temperos, e assim por diante. Na proposta, essa decisão sai de um mini-concurso: dividir os ingredientes da cozinha ao meio, cozinhar com uma metade, provar com a outra, inverter as metades e repetir tudo cinco vezes, aplicando um teste estatístico a cada par de receitas.

A receita escolhida é então feita com todos os ingredientes da cozinha e provada no lote guardado. A diferença de nota entre A e B naquele lote é uma observação. As $KS$ observações alimentam o teste $t$ corrigido.

O ponto crítico está em o que cada nível precisa entregar. O nível interno precisa de uma \emph{escolha}, não de um valor-$p$. O nível externo precisa de uma \emph{estimativa com incerteza honesta}. A Figura~\ref{fig:aninhamento} mostra a estrutura.

\begin{figure}[t]
\centering
\includestandalone[width=\linewidth]{figures/fig_aninhamento}
\caption{Aninhamento proposto. No nível externo, cada repetição $s$ divide os grupos de homologia em $K$ folds. O fold $k$ é o teste externo; os outros formam o treino externo $T_{k,s}$. No nível interno, o 5$\times$2cv divide $T_{k,s}$ ao meio cinco vezes e compara configurações. A configuração escolhida é reajustada em $T_{k,s}$ e avaliada uma única vez no fold $k$, gerando a diferença $x_{k,s}$. As $KS$ diferenças são combinadas pelo teste $t$ corrigido. A seta tracejada marca o único ponto em que o fold de teste é usado.}
\label{fig:aninhamento}
\end{figure}

\subsection{Formalização}

Sejam:
\begin{itemize}[leftmargin=1.4em]
\item $\mathcal{D}$ um dataset com $n$ registros (positivos e controles), organizados em grupos $g\in\mathcal{G}$. No projeto, um grupo é uma componente conexa do grafo de dependência: pares positivo--controle, alinhamentos e sobreposição genômica.
\item A e B dois \emph{procedimentos} de aprendizagem. Um procedimento inclui representação, arquitetura, algoritmo de treino, espaço de busca $\mathcal{C}$ e regra de seleção.
\item Para cada repetição $s=1,\dots,S$, os grupos são permutados com a semente de particionamento $s$ e repartidos em $K$ folds de grupos inteiros, estratificados por classe.
\item Para cada fold $k$: $V_{k,s}$ é o fold de teste externo e $T_{k,s}=\mathcal{D}\setminus V_{k,s}$ é o treino externo, com $n_1=|T_{k,s}|$ e $n_2=|V_{k,s}|$.
\end{itemize}

Para cada célula $(k,s)$ e cada procedimento:
\begin{enumerate}[leftmargin=1.4em]
\item \textbf{Seleção interna.} Escolher $c^\ast_{k,s}\in\mathcal{C}$ usando só $T_{k,s}$. A proposta usa 5$\times$2cv com o $F$ de Alpaydin.
\item \textbf{Reajuste.} Treinar $c^\ast_{k,s}$ em $T_{k,s}$ inteiro.
\item \textbf{Avaliação.} Calcular o MCC em $V_{k,s}$:
\begin{equation}
\MCC=\frac{TP\cdot TN-FP\cdot FN}{\sqrt{(TP+FP)(TP+FN)(TN+FP)(TN+FN)}}.
\end{equation}
\end{enumerate}
A observação externa é a diferença pareada $x_{k,s}=\MCC_A(V_{k,s})-\MCC_B(V_{k,s})$. As duas parcelas são pareadas porque A e B usam o mesmo treino e o mesmo teste.

\subsection{Função de cada nível e necessidade do teste interno}

\paragraph{Nível externo.} Estima o desempenho de cada procedimento e a diferença entre eles. A escolha interna faz parte do procedimento avaliado, e por isso a estimativa não sofre o viés de seleção descrito por \citet{varma2006}.

\paragraph{Nível interno.} Escolhe $c^\ast_{k,s}$. O teste interno \textbf{não é necessário} e, na forma proposta, atrapalha.

\rec\ Há cinco razões.
\begin{enumerate}[leftmargin=1.4em]
\item \textbf{Escolher não é testar.} O resultado do nível interno é uma configuração, não uma afirmação científica. Uma regra como ``a de maior MCC interno médio; em empate, a mais simples'' já basta. O valor-$p$ interno nunca é reportado como evidência, porque a inferência válida é a externa.
\item \textbf{Multiplicidade.} Com $|\mathcal{C}|=C$ configurações, comparar todos os pares exige $C(C-1)/2$ testes $F$ por célula: com $C=8$, são 28 testes em cada uma das $KS$ células. Sem correção, as decisões viram ruído. Com correção, o poder some.
\item \textbf{``Não significativo'' não é ``equivalente''.} Uma regra do tipo ``mantenha a configuração mais simples a menos que a mais complexa seja significativamente melhor'' trata ausência de evidência como evidência de ausência. Isso só é legítimo com uma margem de equivalência declarada \citep{lakens2017}.
\item \textbf{Metade dos dados.} No 5$\times$2cv, cada ajuste interno usa 50\% de $T_{k,s}$, ou seja, 40\% do dataset quando $K=5$. Para redes de cápsulas, isso pode mudar qual configuração é melhor, como o número de épocas ou a regularização, em relação ao reajuste com 100\% de $T_{k,s}$.
\item \textbf{Custo.} São 10 ajustes por configuração, contra 5 num 5-fold interno (\S\ref{sec:custo}).
\end{enumerate}

\subsection{Formação dos conjuntos de treino, validação e teste}
\label{sec:conjuntos}

A Tabela~\ref{tab:conjuntos} mostra o papel de cada conjunto. A regra geral é que nenhuma decisão baseada em dados pode olhar para $V_{k,s}$.

\begin{table}[htbp]
\centering\small
\caption{Papel de cada conjunto no desenho proposto e no protocolo vigente do projeto.}
\label{tab:conjuntos}
\begin{tabularx}{\linewidth}{L{2.5cm}YY}
\toprule
Conjunto & Proposta & Protocolo vigente \proj \\
\midrule
Teste externo & Fold $V_{k,s}$; cada registro é teste $S$ vezes & 20\% dos grupos, fixo, selado até o fim (\texttt{homology\_grouping\_v1}) \\
Treino externo & $T_{k,s}$, com $(K-1)/K$ dos grupos & 80\% dos grupos \\
Treino interno & Metade $A_i$ ou $B_i$ de $T_{k,s}$ & 4 de 5 folds internos por grupos (semente 42) \\
Validação interna & A outra metade & 1 de 5 folds internos \\
Parada antecipada & Dentro de cada ajuste interno, na validação interna; no reajuste, por regra fixa de épocas & Época = mediana inferior das melhores épocas dos 5 folds (máx.\ 50, paciência 15) \\
Limiar & Fixo em 0,5 ou escolhido em predições fora do fold dentro de $T_{k,s}$ & 0,5 primário; limiar de treino só como análise secundária \\
\bottomrule
\end{tabularx}
\end{table}

\subsection{Equações: o teste \texorpdfstring{$t$}{t} 5\texorpdfstring{$\times$}{x}2cv de Dietterich e o \texorpdfstring{$F$}{F} combinado de Alpaydin}
\label{sec:eq5x2}

\lit\ No 5$\times$2cv \citep{dietterich1998}, os dados são permutados cinco vezes ($i=1,\dots,5$) e divididos em duas metades $S_i^{(1)}$ e $S_i^{(2)}$. Na rodada $j=1$, A e B treinam em $S_i^{(1)}$ e são avaliados em $S_i^{(2)}$; na rodada $j=2$, o inverso. A Figura~\ref{fig:reamostragem}a mostra o esquema. Definimos:
\begin{align}
p_i^{(j)} &= \text{diferença de desempenho (A $-$ B) na rodada $j$ da replicação $i$},\\
\bar{p}_i &= \tfrac{1}{2}\big(p_i^{(1)}+p_i^{(2)}\big),\qquad
s_i^2=\big(p_i^{(1)}-\bar{p}_i\big)^2+\big(p_i^{(2)}-\bar{p}_i\big)^2 .
\end{align}
O original usa diferenças de \emph{taxa de erro}; usar diferenças de MCC é uma extensão sem validação nas simulações originais. O termo $s_i^2$ estima a variância da diferença dentro da replicação $i$. Sob $H_0$ (nenhuma diferença esperada), e supondo $p_i^{(j)}\sim\mathcal{N}(0,\sigma^2)$, cada $s_i^2/\sigma^2$ segue aproximadamente uma $\chi^2_1$.

\paragraph{Teste $t$ 5$\times$2cv de Dietterich.}
\begin{equation}
\tilde{t}=\frac{p_1^{(1)}}{\sqrt{\tfrac{1}{5}\sum_{i=1}^{5}s_i^2}}\ \sim\ t_5\quad\text{sob }H_0 .
\end{equation}
O numerador usa \emph{uma única} diferença, a da primeira rodada da primeira replicação. O teste é bilateral e mantém o sinal da diferença. Seu ponto fraco é depender de qual diferença ocupa a posição $(1,1)$, que é arbitrária.

\paragraph{Teste $F$ combinado de Alpaydin.}
\begin{equation}
f=\frac{\sum_{i=1}^{5}\sum_{j=1}^{2}\big(p_i^{(j)}\big)^2}{2\sum_{i=1}^{5}s_i^2}\ \sim\ F_{10,\,5}\quad\text{sob }H_0 .
\end{equation}
O numerador soma os quadrados das dez diferenças: sob $H_0$, $\sum(p_i^{(j)})^2/\sigma^2\approx\chi^2_{10}$. O denominador soma os cinco $s_i^2$: $\sum s_i^2/\sigma^2\approx\chi^2_5$. A razão de duas qui-quadrado divididas por seus graus de liberdade dá $F_{10,5}$. Rejeita-se $H_0$ para valores grandes de $f$.

\lit\ \citet{alpaydin1999} relata erro tipo I menor e poder maior que o teste $t$ de Dietterich, porque o $F$ usa as dez diferenças e não depende da ordem.

\paragraph{Como interpretar.} O $F$ é \emph{não direcional}: ele detecta que as diferenças se afastam de zero, mas não diz quem é melhor. O sinal vem da média das dez diferenças. Os dois testes respondem a uma pergunta sobre \emph{procedimentos} treinados com metade dos dados, num único dataset.

\paragraph{Pressupostos comuns.}
\begin{itemize}[leftmargin=1.4em]
\item As diferenças são aproximadamente normais, com a mesma variância.
\item O numerador e o denominador são aproximadamente independentes, o que não é exato porque as metades se reutilizam.
\item Os registros são permutáveis entre as metades. Com grupos de homologia, a divisão precisa ser por grupos, e as metades deixam de ter tamanhos exatamente iguais.
\end{itemize}

\begin{exemplo}{Dietterich e Alpaydin com os mesmos dados}
\fic\ Suponha as dez diferenças de MCC (A $-$ B) abaixo.
\begin{center}\small
\begin{tabular}{crrrr}
\toprule
$i$ & $p_i^{(1)}$ & $p_i^{(2)}$ & $\bar{p}_i$ & $s_i^2$\\
\midrule
1 & 0,020 & 0,010 & 0,0150 & 0,0000500\\
2 & 0,015 & 0,025 & 0,0200 & 0,0000500\\
3 & 0,030 & 0,005 & 0,0175 & 0,0003125\\
4 & 0,010 & 0,020 & 0,0150 & 0,0000500\\
5 & 0,025 & 0,015 & 0,0200 & 0,0000500\\
\bottomrule
\end{tabular}
\end{center}
O teste $t$ de Dietterich dá $\tilde{t}=1{,}975$ ($p=0{,}105$, 5 g.l.), e o $F$ de Alpaydin dá $f=3{,}537$ ($p=0{,}088$, 10 e 5 g.l.).

Agora troque apenas $p_1^{(1)}$ por 0,002. O $t$ de Dietterich cai para $0{,}201$ ($p=0{,}85$), enquanto o $F$ quase não muda: $3{,}265$ ($p=0{,}10$).

O mesmo conjunto de dados leva a conclusões diferentes no $t$ conforme a ordem das rodadas. Essa instabilidade foi o que motivou o $F$ combinado. Note também que, nos dois casos, as dez diferenças são todas positivas e nenhum teste rejeita $H_0$ a 5\%: com só cinco graus de liberdade no denominador, o poder é baixo.
\end{exemplo}

\begin{figure}[t]
\centering
\includestandalone[width=\linewidth]{figures/fig_5x2_kfold}
\caption{Duas formas de reamostragem. (a) No 5$\times$2cv, cada replicação $i$ é uma nova permutação dividida em metades A e B, usadas alternadamente como treino e teste. (b) No $K$-fold repetido ($K=5$, $S=3$), o número em cada célula é o fold do registro na repetição $s$. Cada fold é teste uma vez por repetição, e cada treino contém $(K-1)/K$ dos dados. Num desenho por grupos, cada quadrado representa um grupo de homologia inteiro.}
\label{fig:reamostragem}
\end{figure}

\subsection{O teste \texorpdfstring{$t$}{t} corrigido compatível com \texorpdfstring{$K$}{K}-fold repetido}
\label{sec:tcorrigido}

\paragraph{Por que corrigir.} Com $J=KS$ diferenças $x_{k,s}$, de média $\bar{x}$ e variância amostral $\hat{\sigma}^2$, o teste $t$ pareado ingênuo usa
\begin{equation}
t_{\text{ingênuo}}=\frac{\bar{x}}{\sqrt{\hat{\sigma}^2/J}},\qquad J-1\ \text{g.l.}
\end{equation}

\lit\ Esse teste supõe as $J$ diferenças independentes, e elas não são: os treinos se sobrepõem, e cada registro aparece em $S$ conjuntos de teste. \citet{bouckaert2004} relatam erro tipo I muito alto para esse teste e recomendam não usá-lo. \citet{bengio2004} provam que não existe estimador não viesado da variância do $K$-fold que valha para toda distribuição.

\paragraph{A correção de Nadeau e Bengio.}
\lit\ Suponha que cada $x_j$ tenha variância $\sigma^2$ e que a correlação entre dois deles seja $\rho$. Então:
\begin{equation}
\Var(\bar{x})=\sigma^2\Big(\rho+\frac{1-\rho}{J}\Big),\qquad \mathbb{E}\big[\hat{\sigma}^2\big]=(1-\rho)\,\sigma^2 .
\end{equation}
Combinando as duas expressões:
\begin{equation}
\widehat{\Var}(\bar{x})=\Big(\frac{1}{J}+\frac{\rho}{1-\rho}\Big)\hat{\sigma}^2 .
\end{equation}
\citet{nadeau2003} propõem a aproximação heurística $\rho\approx n_2/(n_1+n_2)$, ou seja, a fração de teste. Com ela, $\rho/(1-\rho)=n_2/n_1$.

\paragraph{A forma para $K$-fold repetido.} \citet{bouckaert2004} estendem a correção ao $K$-fold repetido $r$ vezes. Com as diferenças $x_{k,s}$ para o fold $k$ da repetição $s$:
\begin{equation}
\boxed{\;t_{\text{corr}}=\frac{\dfrac{1}{KS}\displaystyle\sum_{k=1}^{K}\sum_{s=1}^{S}x_{k,s}}{\sqrt{\Big(\dfrac{1}{KS}+\dfrac{n_2}{n_1}\Big)\hat{\sigma}^2}},\qquad KS-1\ \text{graus de liberdade},\;}
\end{equation}
em que $\hat{\sigma}^2=\frac{1}{KS-1}\sum_{k,s}(x_{k,s}-\bar{x})^2$. No $K$-fold, $n_2/n_1=1/(K-1)$, ou seja, $1/4$ para $K=5$ e $1/9$ para $K=10$.

\lit\ \citet{bouckaert2004} relatam que essa versão tem erro tipo I próximo do nominal e replicabilidade melhor que as alternativas que estudaram, a custo computacional comparável.

\paragraph{O que muda na prática.} O termo $n_2/n_1$ domina o denominador: com $K=5$ e $S=10$, ele vale $0{,}25$, contra $1/50=0{,}02$ do termo ingênuo. Aumentar $S$ reduz só a parcela $1/(KS)$. Por isso o teste corrigido continua conservador, e repetir mais não traz poder ilimitado. O Exemplo da \S\ref{sec:exemplos} mostra o efeito em números.

\paragraph{Pressupostos e limitações que permanecem.}
\begin{enumerate}[leftmargin=1.4em]
\item \textbf{A correlação é aproximada.} A correlação real entre as diferenças depende da estabilidade do algoritmo, e $\rho=n_2/n$ é só uma aproximação \citep{nadeau2003,bengio2004}. O teste é aproximadamente válido, não exato.
\item \textbf{O estimando é o procedimento.} O teste compara A e B como \emph{procedimentos} treinados em conjuntos de tamanho $n_1$ daquela população, não os modelos finais \citep{bates2024}.
\item \textbf{Grupos.} Com dependência entre registros, como pares, homólogos ou regiões sobrepostas, os folds precisam ser de grupos inteiros \citep{roberts2017,teufel2023}. Ainda assim, a correção não modela a correlação extra dentro dos grupos. A razão $n_2/n_1$ passa a ser a razão realizada, e a heurística fica menos justificada quando os grupos têm tamanhos muito desiguais, como o grupo encadeado de 1.154 registros do \emph{L.\ interrogans} (F-120).
\item \textbf{MCC por fold.} O MCC é limitado a $[-1,1]$, não é aditivo e fica ruidoso em folds pequenos. A normalidade de $x_{k,s}$ precisa ser verificada. Uma alternativa é o MCC agregado por repetição: as predições fora do fold de cada $s$ geram $S$ valores, mas isso muda a estrutura de correlação, e a correção acima deixa de ser a mesma.
\item \textbf{Um dataset por vez.} O teste não diz nada sobre outros datasets, organismos ou fontes. A generalização entre datasets exige uma análise por cluster de dependência \citep{demsar2006}.
\item \textbf{Sementes de execução.} Se cada célula tiver $R$ ajustes com sementes diferentes, deve-se fazer a média dentro da célula antes do teste. Usar as $KSR$ diferenças com $J=KSR$ trataria sementes como réplicas independentes.
\end{enumerate}

\rec\ Uma alternativa com a mesma correção é o teste $t$ bayesiano correlacionado \citep{benavoli2017}. Ele produz a probabilidade posterior de que a diferença caia numa região de equivalência prática, o que conversa diretamente com a margem $\delta$ da \S\ref{sec:discussao}.

\subsection{Sementes: particionamento, inicialização e amostragem de controles}
\label{sec:sementes}

Há pelo menos quatro fontes de aleatoriedade, e elas não são intercambiáveis (Tabela~\ref{tab:sementes}).

\begin{table}[htbp]
\centering\small
\caption{Fontes de aleatoriedade, seus papéis no projeto e na proposta.}
\label{tab:sementes}
\begin{tabularx}{\linewidth}{L{2.4cm}YYY}
\toprule
Fonte & O que varia & Projeto hoje \proj & Na proposta \rec \\
\midrule
Particionamento ($s$) & Quais grupos caem em cada fold & Semente 42, partições materializadas e com hash; nunca regeneradas com sementes de execução & $S$ sementes de particionamento distintas, fixadas antes; cada uma gera um manifesto de folds com hash \\
Execução ($r$) & Pesos iniciais, ordem dos lotes, dropout, subamostragem do XGBoost & 10 sementes: 13, 29, 42, 71, 101, 137, 173, 211, 251, 293; reserva só em bloco & As mesmas sementes para todos os modelos de uma comparação; média dentro da célula $(k,s)$ \\
Amostragem de controles ($q$) & Quais janelas genômicas viram controles & Uma semente por procedimento: 20261021 (genômico) e 20261102 (CDS) & Fixa no confirmatório. Variar $q$ cria \emph{outro dataset} e entra só como sensibilidade \\
Reamostragem ($b$) & Índices do bootstrap ou da permutação & 20260715 no protocolo v4 & Inalterada \\
\bottomrule
\end{tabularx}
\end{table}

\rec\ Para separar as fontes de variância, o particionamento e a execução devem ser \emph{cruzados}, ou ao menos ter papéis distintos. Usar a mesma semente para as duas, como em $s=r$, \emph{confunde} os efeitos: não se sabe se uma variação veio da partição ou da inicialização.

\lit\ \citet{bouthillier2021} mostram que a variação devida à amostragem dos dados costuma ser grande perto da variação devida à inicialização. Por isso, variar só as sementes de execução subestima a incerteza de um benchmark.

\proj\ A política de sementes do projeto já proíbe regenerar partições com sementes de execução e trata as sementes como medida de \emph{estabilidade algorítmica}, não como unidades de replicação (\arq{specs/seed_policy.md}).

\subsection{O que é pareado e qual é a unidade de análise}
\label{sec:unidade}

\paragraph{Pareamentos legítimos.}
\begin{itemize}[leftmargin=1.4em]
\item Dentro de cada célula $(k,s)$, A e B usam o mesmo treino e o mesmo teste, então $x_{k,s}$ é pareado.
\item Dentro de uma semente $r$, os modelos estocásticos usam a mesma semente. Um modelo determinístico é ajustado uma vez, e seu resultado é subtraído de cada semente.
\item Dentro de um conjunto de teste fixo, as predições recaem sobre os mesmos registros. É isso que permite o bootstrap pareado.
\end{itemize}

\paragraph{Unidades de análise.}
\begin{itemize}[leftmargin=1.4em]
\item \textbf{Dentro de um dataset}, a unidade de amostragem é o \emph{grupo de homologia} do teste, reamostrado no bootstrap por grupos \citep{field2007}. No desenho proposto, a unidade são as $KS$ células, que são dependentes e por isso exigem a correção.
\item \textbf{Entre datasets}, a unidade é o \emph{cluster de dependência}. \proj\ O manifesto v4 registra 12 clusters entre os 18 datasets elegíveis, e 14 com os candidatos reconstruídos.
\end{itemize}

\paragraph{Como não pseudorreplicar.}
\begin{enumerate}[leftmargin=1.4em]
\item Reduzir as sementes a um efeito por célula ou por dataset antes de qualquer teste.
\item Reduzir os datasets de um mesmo cluster a um efeito por cluster.
\item Nunca juntar folds, sementes e datasets numa única amostra para Wilcoxon ou $t$.
\item Nunca aplicar um teste entre datasets a valores por fold.
\end{enumerate}

\subsection{Hiperparâmetros, parada antecipada e limiar sem tocar no teste externo}

\rec\ Regras que valem tanto para a proposta quanto para o protocolo vigente:
\begin{enumerate}[leftmargin=1.4em]
\item \textbf{Só o treino decide.} Toda operação que aprende com dados fica dentro de $T_{k,s}$ (ou da partição de treino) e é refeita em cada célula externa: normalização, seleção de atributos, calibração, escolha de épocas e de limiar.
\item \textbf{Parada antecipada em dois tempos.} Nos ajustes internos, monitora-se a validação interna. No reajuste, que não tem validação, usa-se uma regra fixa. \proj\ O projeto usa a mediana inferior das melhores épocas dos cinco folds internos.
\item \textbf{Limiar.} O primário é fixo em 0,5. Um limiar alternativo sai de predições fora do fold \emph{dentro} do treino e é reportado como análise secundária.
\item \textbf{Orçamento igual.} O espaço de busca e o número de ajustes por modelo são fixados antes e são iguais, ou explicitamente comparáveis, entre os modelos. \proj\ Hoje o XGBoost e a regressão logística usam hiperparâmetros fixos, e o CapsProm só seleciona a época.
\item \textbf{Construção dos dados vem antes.} A exclusão de controles a menos de 100 nt de qualquer TSS da tabela completa usa as coordenadas de todos os positivos, inclusive os que vão para o teste. Isso é aceitável porque define a população antes da partição e não usa nenhuma predição. O dataset fica congelado antes de qualquer modelo.
\end{enumerate}

%=====================================================================
\section{Perguntas de pesquisa e hipóteses}
\label{sec:hipoteses}
%=====================================================================

A Tabela~\ref{tab:hipoteses} relaciona cada pergunta à hipótese, à comparação e à análise. Há uma distinção essencial:
\begin{itemize}[leftmargin=1.4em]
\item As \textbf{hipóteses de desempenho preditivo} (Q1 a Q8) são sobre modelos num conjunto de teste definido.
\item As \textbf{hipóteses sobre a qualidade dos controles e do pareamento} (Q9) são sobre os \emph{dados}: balanceamento, diversidade e atalhos. Elas se testam sem treinar o modelo de interesse.
\end{itemize}

\begin{landscape}
{\centering\scriptsize
\captionof{table}{Perguntas, hipóteses e análises. $\Delta$ indica diferença pareada de MCC; ``cluster'' é o cluster de dependência do manifesto v4. As linhas sombreadas são perguntas que o desenho proposto, sozinho, não responde.}
\label{tab:hipoteses}
\begin{tabularx}{\linewidth}{L{2.6cm}L{3.6cm}YL{2.1cm}L{1.5cm}L{3.6cm}L{3.0cm}}
\toprule
Pergunta & $H_0$ / $H_1$ & Comparação & Unidade & Métrica & Método & Alcance da conclusão \\
\midrule
Q1. CapsProm supera os baselines? & $H_0$: mediana do $\Delta$ por cluster $=0$; $H_1$: $\neq0$ & CapsProm vs.\ RL, XGBoost (e CNN, se predeclarada), com os mesmos dados, folds, sementes e orçamento & Cluster (entre); grupo de teste (dentro) & MCC; secundárias: AUC, PR-AUC & $\Delta$ por dataset com IC por bootstrap de grupos; Wilcoxon sobre clusters; Holm & Datasets elegíveis, janelas 81/251 nt, controles nativos \\
Q2. Qual variante do CapsProm é melhor? & $H_0$: as variantes têm a mesma distribuição de MCC & Básico vs.\ \emph{Embedding} vs.\ multicabeça & Cluster & MCC & Friedman, depois só os Wilcoxon predeclarados, com Holm & As variantes testadas \\
Q3. Um componente importa (ablação)? & $H_0$: remover o componente não muda o MCC; para ``pode remover'': $H_0$: $|\Delta|\geq\delta$ & Modelo completo vs.\ sem roteamento dinâmico, sem decodificador ou com outro número de cápsulas; um fator por vez & Cluster & MCC & Wilcoxon e Hodges--Lehmann com IC; família própria de Holm; TOST com margem $\delta$ & O componente nesta arquitetura e neste treino \\
Q4. Estabilidade à inicialização & Descritiva; ``não colapsa'' é uma afirmação de limite & 10 sementes de execução, com dados e partição fixos & Semente dentro do dataset (descritivo) & DP e IQR do MCC; taxa de colapso & IQR; limite superior exato $1-0{,}05^{1/n}$ (26\% com $n=10$) & Esses datasets e sementes \\
\rowcolor{black!6}
Q5. Estabilidade ao particionamento & Descritiva: variância do $\Delta$ entre partições & $K$-fold por grupos repetido \emph{dentro do treino} & Repetição $s$ (dependente) & MCC & Componentes de variância; $t$ corrigido só como descrição & O procedimento, no tamanho $n_1$ \\
\rowcolor{black!6}
Q6. Generaliza entre datasets, fontes e organismos? & $H_0$: a perda de desempenho fora do domínio é igual para A e B & Treinar em $a$ e testar em $b$; ou deixar um cluster de fora & Par fonte$\to$alvo, agregado por cluster-alvo & MCC & Matriz de transferência; Wilcoxon se houver $\geq10$ clusters-alvo & Pares de domínios incluídos \\
\rowcolor{black!6}
Q7. A estratégia de controles de treino muda o desempenho? & $H_0$: $\Delta^{\text{est}}=0$ num painel de teste comum & Treino com controles nativos vs.\ genômicos GC vs.\ multinegativo, com modelo e treino fixos & Grupo de teste; cluster & MCC por painel & Bootstrap pareado no painel comum; agregação por cluster & Só datasets com as duas estratégias (hoje 3 bactérias): exploratória \\
\rowcolor{black!6}
Q8. Há interação arquitetura $\times$ controles? & $H_0$: $\Delta^{\text{mod}}_G-\Delta^{\text{mod}}_N=0$ & Desenho fatorial $2\times2$ (ou maior) no painel comum & Cluster & MCC & Contraste diferença de diferenças com IC por bootstrap & Exploratória (protocolo v1 de controles) \\
\rowcolor{black!6}
Q9. Os controles são bem pareados e livres de atalhos? & $H_0$ de equivalência: $|\text{SMD}|\geq\tau$ (desbalanceado); atalho: MCC do classificador só-GC $\leq\epsilon$ & Controles vs.\ positivos, sem o modelo de interesse & Par ou registro, dentro do dataset & SMD de GC, JSD de $k$-mers, TV posicional, MCC do atalho & Testes de equivalência; IC do MCC do atalho; diagnósticos da \S\ref{sec:diagnosticos} & Propriedades medidas; não prova que os controles são não promotores \\
\bottomrule
\end{tabularx}
}
\end{landscape}

\paragraph{O que o desenho proposto não consegue responder.} \rec\ O $K$-fold externo, com 5$\times$2cv interno, compara procedimentos \emph{dentro} de um dataset. Por isso ele não responde a:
\begin{itemize}[leftmargin=1.4em]
\item \textbf{Generalização entre datasets ou organismos (Q6).} Todos os folds vêm da mesma população.
\item \textbf{Superioridade geral do CapsProm (Q1).} Isso exige agregar datasets por cluster. Somar valores-$p$ ou contar vitórias por dataset não basta \citep{demsar2006}.
\item \textbf{Equivalência.} Um teste de diferença sem margem $\delta$ não sustenta ``A e B são equivalentes'' nem ``o componente pode ser removido''.
\item \textbf{Qualidade e validade biológica dos controles (Q9).} Nenhum valor de MCC prova que um controle é um não promotor. Os controles genômicos têm incerteza do tipo positivo-não-rotulado \citep{elkan2008}: estar ausente da tabela de TSSs não prova inatividade.
\item \textbf{Desempenho numa varredura genômica real.} Falsos positivos por megabase exigem um painel de varredura, previsto no protocolo de controles como análise adicional.
\item \textbf{Mecanismo.} Uma ablação mostra que retirar um componente muda o MCC naquele treino, mas não que o componente ``aprende'' um motivo biológico específico.
\end{itemize}

%=====================================================================
\section{Separar os efeitos da arquitetura e da construção dos dados}
\label{sec:efeitos}
%=====================================================================

\subsection{Critérios de pareamento efetivamente adotados}
\label{sec:criterios}

\proj\ Os critérios abaixo foram conferidos nos contratos e nas especificações.

\paragraph{Controle genômico pareado por GC.} Fonte: \arq{configs/graphpro_reconstruction_v3.yaml} e \arq{specs/negative_control_protocol.md}.
\begin{itemize}[leftmargin=1.4em]
\item Um controle por TSS positivo, na mesma montagem e na mesma fita, com a mesma janela (81 nt: 60 a montante, o TSS e 20 a jusante).
\item Exclusão de qualquer posição a menos de 100 nt de qualquer TSS da tabela primária completa, em todas as classes e nas duas fitas.
\item Sem sobreposição com outro controle e só DNA estrito (ACGT).
\item Sem identidade com um positivo ou com outro controle, em qualquer orientação. Também são rejeitados controles iguais a positivos ativos de outros 18 datasets.
\item Diferença de no máximo 4 bases G/C em relação ao positivo.
\item Amostragem com a semente 20261021: até 100.000 sondagens aleatórias e, depois disso, uma varredura exaustiva com amostragem de reservatório.
\item O par positivo--controle é indivisível em toda filtragem e partição.
\end{itemize}

\paragraph{Controle derivado de CDS.} Fonte: \arq{specs/cds_negative_control_protocol.md}.
\begin{itemize}[leftmargin=1.4em]
\item Janela de 81 nt na fita senso de um CDS elegível, com fase em ciclo 0, 1, 2.
\item Pelo menos 60 nt de distância do início e do fim do CDS.
\item Sem sobreposição com outro CDS ou RNA anotado.
\item As mesmas regras de exclusão de TSS e de GC do controle genômico.
\item O \emph{L.\ interrogans} foi excluído desta condição por esgotar o espaço codificante (F-123).
\end{itemize}

\paragraph{Diagnósticos já medidos.} Apesar da tolerância por par, os controles são mais ricos em GC que os positivos em \emph{C.\ jejuni} e \emph{C.\ pneumoniae} (F-124): 2,3 e 1,9 pontos percentuais no painel genômico, e 2,8 e 2,6 no painel CDS. No \emph{L.\ interrogans} não há diferença.

\paragraph{Agrupamento.} Fonte: \arq{specs/homology_grouping_protocol.md}.
\begin{itemize}[leftmargin=1.4em]
\item As arestas do grafo vêm de três fontes: pares; alinhamentos MMseqs2 com identidade e cobertura de pelo menos 0,80 nas duas fitas, ou distância de Hamming de até $\lfloor0{,}2L\rfloor$; e sobreposição genômica \citep{steinegger2017}.
\item As componentes conexas nunca são divididas entre treino e teste nem entre folds internos \citep{teufel2023,petti2022}.
\end{itemize}

\paragraph{Estatística entre painéis.}
\begin{itemize}[leftmargin=1.4em]
\item Os painéis de controles nunca são concatenados num único MCC.
\item As comparações entre painéis ficam restritas à interseção dos positivos de teste, e essa interseção é fixada antes de qualquer avaliação.
\item A interação modelo $\times$ política é exploratória na versão 1.
\end{itemize}

\subsection{Desenho fatorial}

\rec\ Considere três fatores:
\begin{itemize}[leftmargin=1.4em]
\item a arquitetura $a$ (CapsProm, baselines);
\item a estratégia de controles \emph{de treino} $n$ (nativa, genômica GC, CDS, multinegativa);
\item o painel \emph{de teste} $P$.
\end{itemize}
Seja $\mu_{a,n}(P)$ o MCC esperado no painel $P$. Os três efeitos se definem num painel fixo $P^\ast$ (Figura~\ref{fig:fatorial}):
\begin{align}
\text{efeito do modelo (com $n$ fixo):}\quad & \Delta^{\text{mod}}_n=\mu_{C,n}(P^\ast)-\mu_{B,n}(P^\ast),\\
\text{efeito da estratégia (com $a$ fixo):}\quad & \Delta^{\text{est}}_a=\mu_{a,G}(P^\ast)-\mu_{a,N}(P^\ast),\\
\text{interação:}\quad & I=\Delta^{\text{mod}}_G-\Delta^{\text{mod}}_N=\Delta^{\text{est}}_C-\Delta^{\text{est}}_B .
\end{align}

\begin{enumerate}[leftmargin=1.4em]
\item \textbf{Efeito do modelo.} Fixe os dados, os controles de treino, os folds, as sementes, o orçamento e a regra de seleção. Mude só a arquitetura.
\item \textbf{Efeito dos controles.} Fixe a arquitetura e todo o procedimento de modelagem. Mude só os controles \emph{de treino}, e avalie todos os modelos resultantes no \emph{mesmo} painel de teste. Faça isso em mais de um painel, como o nativo e o genômico, para ver se a conclusão depende do painel.
\item \textbf{Interação.} Use o contraste $I$, com um IC por bootstrap pareado de grupos no painel comum, recalculando as quatro células em cada réplica. Entre datasets, agregue $I$ por cluster. Com poucos clusters, como as três bactérias reconstruídas hoje, a análise é exploratória.
\end{enumerate}

\begin{figure}[t]
\centering
\includestandalone[width=\linewidth]{figures/fig_fatorial}
\caption{Separação de efeitos. (a) Desenho $2\times2$: arquitetura (CapsProm $C$ ou baseline $B$) cruzada com os controles usados no treino (nativos $N$ ou genômicos pareados por GC $G$). Todas as células são avaliadas no mesmo painel de teste $P^\ast$. As setas indicam o efeito da estratégia, dentro de uma arquitetura, e o efeito do modelo, dentro de uma estratégia. A interação é a diferença entre os efeitos do modelo nas duas estratégias. (b) Gráfico de interação com valores fictícios: linhas não paralelas indicam interação.}
\label{fig:fatorial}
\end{figure}

\subsection{Pares e grupos no particionamento}

\rec\ Para manter pares e grupos juntos (Figura~\ref{fig:pares}):
\begin{enumerate}[leftmargin=1.4em]
\item \textbf{Grupo global.} Antes de qualquer partição, construa o grafo com todos os registros de todas as condições: positivos e controles nativos, genômicos e CDS. Isso garante que um positivo e qualquer controle derivado dele fiquem no mesmo lado em \emph{todas} as condições. \proj\ É a regra de ``grupo global'' do protocolo de controles; o controle CDS herda o grupo do positivo.
\item \textbf{Folds de grupos.} Use \emph{StratifiedGroupKFold} sobre os grupos, tanto no nível externo da proposta quanto no interno. No 5$\times$2cv, as metades também têm de ser de grupos inteiros.
\item \textbf{Sobreposição de loci.} Janelas que se sobrepõem no genoma, em qualquer fita, entram como aresta. Um controle que se liga a um registro de outra partição é excluído do painel e registrado, sem mudar o split.
\item \textbf{Grupos grandes.} Um grupo encadeado grande desequilibra os folds. \proj\ É o caso do grupo de 1.154 registros do \emph{L.\ interrogans}, que ficou inteiro no fold interno 4 (F-120). Um desequilíbrio assim aumenta a variância entre folds e merece um relatório à parte; ele é tolerado pelo limite de 20\%.
\end{enumerate}

\begin{figure}[t]
\centering
\includestandalone[width=\linewidth]{figures/fig_pares}
\caption{Preservação de pares e grupos. (a) Cada positivo (círculo) se liga ao seu controle (quadrado). Alinhamentos e sobreposições genômicas unem pares em componentes conexas G1--G4. (b) No particionamento correto, cada componente vai inteira para um único bloco. (c) Na divisão por registro, o par P2--C2 e os homólogos P2--P3 ficam em lados opostos, e o modelo pode ``reconhecer'' o teste por semelhança em vez de generalizar.}
\label{fig:pares}
\end{figure}

\subsection{Positivos comuns e conjuntos de avaliação comuns}

Quando a estratégia de controles muda, o conjunto de \emph{teste} também pode mudar. Há três situações.
\begin{enumerate}[leftmargin=1.4em]
\item \textbf{Mesmo painel para todos.} Os modelos treinados com controles diferentes são avaliados nos mesmos registros. O pareamento é completo e o bootstrap pareado vale. \rec\ Este é o desenho preferido para Q7 e Q8.
\item \textbf{Painéis diferentes com os mesmos positivos.} Por exemplo, o painel nativo e o genômico. A sensibilidade sobre os positivos comuns é pareada, mas a especificidade não, porque os controles são outros. O MCC mistura as duas, então uma diferença de MCC entre painéis reflete a \emph{dificuldade do painel} tanto quanto o modelo. Ela não mede a qualidade dos controles.
\item \textbf{Positivos diferentes.} Acontece quando um positivo perde o controle num painel, como as falhas do painel CDS. \proj\ O protocolo restringe a comparação à interseção dos positivos de teste, com os identificadores fixados pelos registros de construção antes de qualquer avaliação.
\end{enumerate}

\rec\ Para um IC de uma diferença entre painéis, use um \emph{bootstrap conjunto de grupos}: cada unidade reamostrada é um ``supergrupo'' com o positivo, seus controles em todos os painéis e os homólogos ligados a eles. Assim, cada réplica recalcula os dois painéis sobre o mesmo conjunto de positivos. O pareamento é parcial, e o relatório deve dizer isso.

\subsection{Por que um MCC maior não indica controle melhor, e o que medir no lugar}
\label{sec:diagnosticos}

\lit\ Controles fáceis, como sequências de composição diferente, produzem MCC alto porque o modelo aprende um \emph{atalho} \citep{geirhos2020}. \citet{krutzfeldt2020} e \citet{gonzalez2025} mostram que a escolha dos negativos muda o desempenho e o viés dos preditores de sequências regulatórias. Em particular, negativos tirados de CDS induzem um viés de GC em preditores de promotores bacterianos.

\proj\ O protocolo do projeto afirma o mesmo: um MCC maior sob controles mais fáceis não é evidência de melhor detecção de promotores. A Tabela~\ref{tab:diagnosticos} reúne os diagnósticos recomendados.

\begin{table}[htbp]
\centering\small
\caption{Diagnósticos de qualidade dos controles e do pareamento. Nenhum depende do desempenho do modelo de interesse.}
\label{tab:diagnosticos}
\begin{tabularx}{\linewidth}{L{2.7cm}YL{3.4cm}L{3.0cm}}
\toprule
Diagnóstico & Como medir & Sinal de alerta \rec & Situação no projeto \proj \\
\midrule
Balanceamento de GC & Diferença padronizada de médias (SMD) e distribuição de GC por classe, global e no teste & $|\text{SMD}|>0{,}1$ ou diferença sistemática apesar da tolerância & Diferença de 1,9--2,8 p.p.\ em Cje e Cpne (F-124) \\
Composição de $k$-mers & Divergência de Jensen--Shannon entre classes, para $k=2,3$ & Divergência acima de um limiar predeclarado & Auditoria de distribuições (R-051) \\
Perfil posicional & Distância de variação total (TV) por posição & TV alta \emph{fora} das posições esperadas de elementos promotores & TV maior em $-8$, $-35/-34$ e 0, posições esperadas (F-124) \\
Periodicidade de 3 nt & Espectro por classe (painel CDS) & Pico só nos controles & Exigido pelo protocolo CDS \\
Atalho trivial & MCC de um classificador só-GC ou só-$k$-mer, com IC & MCC do atalho próximo ao do modelo & Referência só-GC congelada (R-052) \\
Diversidade e cobertura & Número de grupos, maior grupo, distribuição dos controles no cromossomo & Grupos gigantes; controles concentrados & Grupo encadeado (F-120); viés posicional no CDS do Lint (F-123) \\
Ruído de rótulo & TSSs não anotados entre os controles; condição experimental dos TSSs & Controles perto de TSSs de outra condição & Exclusão pelas duas temperaturas no Lint \\
Dependência externa & Identidade ou alinhamento com outros datasets & Controles iguais a positivos de outro dataset & Manifesto de dependência v4 (F-125) \\
\bottomrule
\end{tabularx}
\end{table}

\rec\ Um bom controle é aquele que (i) se parece com os positivos nas características que \emph{não} definem promotor, como GC, comprimento, fita e região do genoma; (ii) não deixa um atalho trivial; (iii) é diverso e cobre o genoma; e (iv) tem baixa probabilidade de ser um promotor não anotado.

Esses critérios são verificados \emph{antes} e \emph{fora} da comparação de modelos. Um modelo pode ter MCC menor com controles melhores, e é exatamente isso que se espera quando o atalho some.

%=====================================================================
\section{Vantagens, limitações e alternativas}
\label{sec:alternativas}
%=====================================================================

A Tabela~\ref{tab:alternativas} compara as alternativas. Os custos usam a notação da \S\ref{sec:custo}.

\begin{landscape}
{\centering\scriptsize
\captionof{table}{Alternativas de desenho. ``Procedimento'' inclui a seleção de hiperparâmetros. Os custos são o número de ajustes por modelo, por dataset e por semente de execução; $C$ é o número de configurações e $F$ o número de folds internos.}
\label{tab:alternativas}
\begin{tabularx}{\linewidth}{L{2.9cm}YYL{2.4cm}YY}
\toprule
Alternativa & Prós & Contras & Custo & Pressupostos & Quando usar \\
\midrule
(a) Proposta: $K$-fold repetido externo com $t$ corrigido; 5$\times$2cv interno com $F$ de Alpaydin & Estima o procedimento inteiro, sem viés de seleção; usa todos os dados como teste ao longo das repetições & Teste interno desnecessário, com multiplicidade e só 50\% do treino; correção heurística; um dataset por vez; recicla o teste externo selado & $SK(10C+1)$ & Grupos permutáveis; $\rho\approx n_2/n$; normalidade de $x_{k,s}$ & Datasets pequenos sem teste externo, com custo de treino baixo; não no confirmatório do CapsProm \\
(b) CV aninhada com seleção interna sem teste de hipótese & Separa escolha e estimação \citep{varma2006,cawley2010}; seleção barata e estável & Inferência pelo $t$ corrigido ainda é aproximada; custo alto para redes & $SK(FC+1)$ & Iguais a (a), sem o teste interno & Estimar o desempenho do procedimento quando não há teste externo; estabilidade dentro do treino (S1) \\
(c) 5$\times$2cv como experimento dedicado a dois procedimentos & Simples; erro tipo I controlado para um par \citep{dietterich1998,alpaydin1999} & Treino com 50\% dos dados; poder baixo (5 g.l.); um par e um dataset por vez; feito para taxa de erro & $10$ por procedimento (sem seleção) & Diferenças normais; metades por grupos & Uma pergunta pontual sobre dois procedimentos fixos, com treino barato \\
(d) Teste externo fixo por grupos, com IC por bootstrap de grupos (protocolo vigente) & Teste aberto uma vez; comparável com partições publicadas; barato; IC respeita grupos \citep{field2007} & Incerteza condicionada a uma partição; não mede a variação devida ao particionamento & $FC+1$ & Grupos independentes entre si; teste representativo & Comparação confirmatória por dataset; base da agregação entre datasets \\
(e) Múltiplos modelos e datasets: Wilcoxon e Friedman sobre clusters, com Holm; alternativa bayesiana & Inferência sobre a população de datasets; não paramétrico; controla multiplicidade \citep{demsar2006,holm1979,benavoli2017} & Precisa de clusters suficientes (o projeto exige $\geq10$); poder limitado com 12 & Sem custo de treino adicional & Clusters independentes; efeito simétrico para Wilcoxon & Conclusões gerais (Q1--Q3) \\
\bottomrule
\end{tabularx}
}
\end{landscape}

\subsection{Pontos de discussão}
\label{sec:discussao}

\paragraph{Dependência entre avaliações.}
\lit\ Folds de um mesmo $K$-fold compartilham treino. Repetições compartilham tudo. Registros homólogos são correlacionados, e datasets de uma mesma fonte também. Cada nível de dependência pede seu tratamento:
\begin{itemize}[leftmargin=1.4em]
\item correção de variância para folds e repetições \citep{nadeau2003,bouckaert2004};
\item grupos para registros \citep{teufel2023,roberts2017};
\item clusters para datasets.
\end{itemize}

\paragraph{Viés de seleção.}
\lit\ Escolher a melhor entre muitas configurações e reportar o desempenho dela no mesmo dado infla o resultado \citep{varma2006,cawley2010,krstajic2014}. O aninhamento corrige isso no nível externo. O teste interno não acrescenta proteção.

\paragraph{Poder estatístico.}
O $F$ de Alpaydin tem 5 graus de liberdade no denominador. O $t$ corrigido é conservador por construção. Entre datasets, com 12 clusters, o menor valor-$p$ bilateral possível no Wilcoxon exato é $2/2^{12}\approx0{,}00049$, o que deixa espaço para a correção de Holm com poucas hipóteses, mas pouco poder para efeitos pequenos e inconsistentes.

\rec\ O caminho para mais poder está em mais datasets independentes e em efeitos maiores, não em mais folds.

\paragraph{Comparações múltiplas.}
\proj\ O protocolo aplica Holm à família de pares de modelos predeclarados na métrica primária. As métricas secundárias são exploratórias.

\rec\ As ablações devem formar uma família própria, pequena e predeclarada. O procedimento de Holm \citep{holm1979} rejeita $H_{(i)}$ enquanto $p_{(i)}\leq\alpha/(m-i+1)$, com os valores-$p$ em ordem crescente.

\paragraph{Magnitude do efeito, intervalos e relevância prática.}
\rec\ Reporte sempre $\Delta$MCC com intervalo de confiança, e declare antes uma margem mínima de relevância $\delta$. O valor de $\delta$ é uma decisão científica que o autor precisa justificar: este relatório usa $\delta=0{,}02$ só como exemplo.

\proj\ O erro padrão da média de dez sementes ficou entre 0,002 e 0,0026 de MCC nas rodadas concluídas (\arq{specs/seed_policy.md}). Uma margem dessa ordem de grandeza seria dominada pela variação entre partições e entre datasets, não pela variação entre sementes.

Um intervalo inteiramente dentro de $(-\delta,\delta)$ sustenta equivalência prática; um teste de equivalência (TOST) formaliza isso \citep{lakens2017}. A Figura~\ref{fig:intervalos} mostra os casos possíveis.

\subsection{Número de treinamentos}
\label{sec:custo}

Notação:
\begin{itemize}[leftmargin=1.4em]
\item $K$: folds externos; $S$: repetições externas;
\item $F$: folds internos ($F=5$ no projeto); o 5$\times$2cv equivale a 10 ajustes;
\item $C$: configurações candidatas; $R$: sementes de execução por célula;
\item $M$: modelos; $D$: datasets.
\end{itemize}

Cada célula externa faz a seleção interna e mais um reajuste. Então:
\begin{align}
N_{\text{proposta}} &= M\,D\,S\,K\,R\,(10\,C+1),\\
N_{\text{aninhada}} &= M\,D\,S\,K\,R\,(F\,C+1),\\
N_{\text{teste fixo}} &= M\,D\,R\,(F\,C+1),\\
N_{\text{5}\times\text{2 dedicado}} &= 2\,D\,R\cdot10\quad\text{(dois procedimentos fixos)}.
\end{align}

\begin{table}[htbp]
\centering\small
\caption{Número de treinamentos com $K=5$, $S=10$, $C=8$, $F=5$, $M=4$ e $D=18$. O CapsProm vigente seleciona só a época ($C=1$). Os modelos determinísticos precisariam de $R=1$, o que reduz os totais da coluna $R=10$.}
\label{tab:custo}
\begin{tabularx}{\linewidth}{Y>{\raggedleft\arraybackslash}p{2.3cm}rr}
\toprule
Desenho & Por modelo e dataset ($R=1$) & Total, $R=1$ & Total, $R=10$ \\
\midrule
(a) Proposta & 4.050 & 291.600 & 2.916.000 \\
(b) Aninhada, sem teste interno & 2.050 & 147.600 & 1.476.000 \\
(d) Teste fixo com seleção interna ($C=8$) & 41 & 2.952 & 29.520 \\
(d) Teste fixo, como o CapsProm vigente ($C=1$) & 6 & 432 & 4.320 \\
S1: estabilidade dentro do treino ($S'=3$, $K=5$, configuração fixa) & 15 & 1.080 & 10.800 \\
(c) 5$\times$2cv dedicado (dois procedimentos) & 10 & 360 & 3.600 \\
\bottomrule
\end{tabularx}
\end{table}

A proposta custa cerca de \textbf{2 vezes} a CV aninhada com 5-fold interno e cerca de \textbf{675 vezes} o protocolo vigente do CapsProm ($4.050/6$). \proj\ Para comparação, o plano canônico já soma 1.560 ajustes ($26\times10\times6$).

%=====================================================================
\section{Exemplos fictícios}
\label{sec:exemplos}
%=====================================================================

\fic\ Todos os números desta seção foram inventados para ilustrar a interpretação. Eles não vêm de nenhuma execução do projeto.

\begin{exemplo}{CapsProm vs.\ baseline num dataset}
\textbf{(a) Teste externo fixo, como no protocolo vigente.} A mediana das dez diferenças pareadas por semente é $\Delta\MCC=+0{,}021$. O IC 95\% do bootstrap de grupos é $[-0{,}004;\ 0{,}046]$.
\begin{itemize}[leftmargin=1.2em]
\item O IC contém zero, então a diferença não é significativa a 5\%.
\item O IC também ultrapassa $\delta=0{,}02$, então não é possível afirmar equivalência. A conclusão correta é ``\emph{inconclusivo}'', não ``empate''.
\end{itemize}

\textbf{(b) $K$-fold repetido.} Com $K=5$, $S=10$, média $\bar{x}=0{,}015$ e desvio $\hat{\sigma}=0{,}030$ nas 50 diferenças:
\begin{align*}
t_{\text{ingênuo}}&=\frac{0{,}015}{0{,}030/\sqrt{50}}=3{,}54\quad(p=0{,}0009),\\
t_{\text{corr}}&=\frac{0{,}015}{\sqrt{(1/50+1/4)\cdot0{,}030^2}}=0{,}96\quad(p=0{,}34,\ 49\ \text{g.l.}).
\end{align*}
O teste ingênuo ``acha'' uma diferença que a correção não sustenta. Com esses $\hat{\sigma}$ e $K$, a média precisaria passar de cerca de $0{,}031$ para ser significativa, por maior que fosse $S$.
\end{exemplo}

\begin{exemplo}{Ablação entre variantes do CapsProm}
A ablação compara o CapsProm completo com uma variante sem roteamento dinâmico (uma única iteração). O efeito por cluster de dependência é $\Delta=$ completo $-$ variante, nos 12 clusters:
\[
0{,}021;\ -0{,}005;\ 0{,}012;\ 0{,}004;\ -0{,}011;\ 0{,}018;\ 0{,}009;\ -0{,}003;\ 0{,}015;\ -0{,}008;\ 0{,}006;\ 0{,}010.
\]
O resultado é:
\begin{itemize}[leftmargin=1.2em]
\item mediana $0{,}0075$ e estimador de Hodges--Lehmann $0{,}006$, com IC 95\% $[-0{,}001;\ 0{,}0135]$;
\item 8 vitórias e 4 derrotas;
\item Wilcoxon $W=17$, $p=0{,}092$.
\end{itemize}

Suponha agora uma família predeclarada de três ablações, com valores-$p$ 0,004, 0,030 e 0,200. Ajustados por Holm, eles viram 0,012, 0,060 e 0,200, e só a primeira ablação sustenta efeito.

A ablação do roteamento é \emph{não significativa}, mas isso \emph{não} autoriza removê-lo. Para remover, seria preciso mostrar equivalência: o IC teria de ficar dentro de $(-\delta,\delta)$. Com $\delta=0{,}02$, o IC $[-0{,}001;\ 0{,}0135]$ está dentro, e a equivalência ficaria sustentada \emph{se} $\delta$ tivesse sido declarado antes.
\end{exemplo}

\begin{exemplo}{Estratégias de controles negativos}
O mesmo XGBoost é treinado com controles nativos ou genômicos pareados por GC e avaliado nos dois painéis:
\begin{center}\small
\begin{tabular}{lcc}
\toprule
MCC & painel nativo & painel genômico GC \\
\midrule
Treino com controles nativos & 0,78 & 0,41 \\
Treino com controles genômicos GC & 0,70 & 0,62 \\
Classificador só-GC (referência de atalho) & 0,46 & 0,04 \\
\bottomrule
\end{tabular}
\end{center}
\begin{itemize}[leftmargin=1.2em]
\item \textbf{Leitura correta.} No painel genômico comum, o treino genômico supera o nativo em $0{,}21$; essa comparação é pareada. O classificador só-GC atinge 0,46 no painel nativo, de modo que boa parte do 0,78 pode vir de composição, e não de sinal promotor.
\item \textbf{Leitura incorreta.} Comparar 0,78 com 0,62 e concluir que os controles nativos são ``melhores''. Os painéis têm controles diferentes; o número maior indica um painel mais fácil.
\end{itemize}
\end{exemplo}

\begin{exemplo}{Significância, relevância e equivalência}
A Figura~\ref{fig:intervalos} mostra quatro intervalos para a mesma margem $\delta=0{,}02$:
\begin{enumerate}[label=(\roman*), leftmargin=1.6em]
\item Diferença significativa e relevante.
\item Diferença significativa, mas com o IC inteiro dentro de $\pm\delta$. É praticamente irrelevante: há equivalência prática \emph{e} uma diferença detectável.
\item Diferença não significativa e inconclusiva: o IC é largo demais para dizer qualquer coisa.
\item Diferença não significativa e equivalente: o IC é estreito e está dentro de $\pm\delta$.
\end{enumerate}
Só o caso (iv) autoriza ``os procedimentos são equivalentes'', e apenas porque a margem foi declarada antes.
\end{exemplo}

\begin{figure}[h]
\centering
\includestandalone[width=0.92\linewidth]{figures/fig_intervalos}
\caption{Significância não é relevância, e ausência de significância não é equivalência. Cada barra é um IC 95\% fictício de $\Delta$MCC. A faixa verde é a região de equivalência $\pm\delta$, que precisa ser declarada antes da análise. Um intervalo que cruza zero mas sai da faixa, como em (iii), é inconclusivo.}
\label{fig:intervalos}
\end{figure}

%=====================================================================
\section{Recomendação: um protocolo viável}
\label{sec:recomendacao}
%=====================================================================

\rec\ A Tabela~\ref{tab:recomendacao} prioriza os experimentos. A Figura~\ref{fig:protocolo} mostra o fluxo. A recomendação central é manter a linha confirmatória do protocolo v4 e tratar as peças úteis da proposta como análises secundárias, que nunca alimentam escolhas confirmatórias.

\begin{figure}[t]
\centering
\includestandalone[width=\linewidth]{figures/fig_protocolo}
\caption{Protocolo recomendado. À esquerda, a linha confirmatória, que coincide com o protocolo v4 do projeto, com a seleção interna feita por regra e sem teste de hipótese. À direita, as análises secundárias que aproveitam partes da proposta. S1 e S2 usam só a partição de treino e, portanto, não reciclam o teste externo selado.}
\label{fig:protocolo}
\end{figure}

\begin{table}[htbp]
\centering\small
\caption{Recomendação priorizada. ``C'' indica análise confirmatória; ``E'', exploratória; ``D'', descritiva. Uma análise só é confirmatória se estiver predeclarada num contrato versionado antes de abrir o teste.}
\label{tab:recomendacao}
\begin{tabularx}{\linewidth}{cL{2.7cm}YYc}
\toprule
\# & Pergunta & Experimento & Análise & Tipo \\
\midrule
1 & Q1: CapsProm vs.\ baselines & Protocolo v4: teste externo por grupos, seleção interna 5-fold por grupos, 10 sementes, mesmo orçamento & $\Delta$MCC com IC por bootstrap de grupos; Wilcoxon sobre clusters; Holm & C \\
2 & Q3: ablações (no máximo 3 ou 4, predeclaradas) & Um fator por vez, mesmos dados, folds, sementes e orçamento & Igual a 1, com família de Holm própria; TOST com margem $\delta$ para afirmar remoção & C \\
3 & Q9: qualidade dos controles & Diagnósticos sem modelo e classificador de atalho (Tabela~\ref{tab:diagnosticos}) & SMD, JSD, TV e MCC do atalho com IC, com limiares de alerta fixados antes & D \\
4 & Q4: estabilidade à inicialização & As 10 sementes já previstas & IQR, DP e limite exato de colapso & D \\
5 & Q7 e Q8: controles e interação & Fatorial com painel comum nas bactérias reconstruídas & Bootstrap pareado no painel comum; contraste de interação & E \\
6 & Q5: estabilidade ao particionamento & S1: $K$-fold por grupos repetido \emph{dentro do treino} ($S'=3$, sementes de particionamento próprias) & Componentes de variância; $t$ corrigido ou $t$ bayesiano correlacionado como descrição & E \\
7 & Q6: generalização entre organismos & Treino em um dataset e teste em outro, com a mesma janela; novo contrato & Matriz de transferência; Wilcoxon se houver $\geq10$ clusters-alvo & E \\
8 & Par pontual de procedimentos & S2: 5$\times$2cv com $F$ de Alpaydin, só se a pergunta for de fato sobre um par e o custo permitir & $F_{10,5}$ e média das diferenças & E \\
\bottomrule
\end{tabularx}
\end{table}

\subsection{Decisões a fixar antes de qualquer execução}

\rec\ Cada item abaixo deve constar de um contrato YAML versionado. \proj\ Pelas regras do repositório, mudar qualquer um deles depois exige uma nova versão de configuração e um novo caminho de estado.

\begin{enumerate}[leftmargin=1.6em]
\item \textbf{Escopo:} datasets elegíveis, clusters de dependência e critérios de exclusão.
\item \textbf{Família confirmatória:} pares de modelos, ablações predeclaradas, métrica primária (MCC), limiar de 0,5 e correção de Holm.
\item \textbf{Seleção:} espaço de busca $\mathcal{C}$, orçamento por modelo, regra de seleção sem teste de hipótese com critério de desempate, e regra de épocas.
\item \textbf{Sementes:} particionamento, execução, amostragem de controles e reamostragem, com papéis distintos; reserva usada só em bloco.
\item \textbf{Relevância:} margem $\delta$ e se haverá testes de equivalência, com a justificativa científica.
\item \textbf{Incerteza:} unidade do bootstrap (grupo de homologia), número de réplicas e tratamento de réplicas sem uma das classes.
\item \textbf{Painéis:} quais painéis de teste existem, a interseção de positivos para comparações entre painéis e a proibição de concatená-los.
\item \textbf{Controles:} diagnósticos e limiares de alerta, e o classificador de atalho de referência.
\item \textbf{Falhas:} sementes que falham são reportadas e nunca substituídas; casos incompletos ficam fora do par primário.
\item \textbf{Fronteira:} o que é confirmatório e o que é exploratório, por escrito, antes de abrir o teste.
\end{enumerate}

\subsection{O que não fazer}
\begin{itemize}[leftmargin=1.4em]
\item Escolher hiperparâmetros com teste de hipótese no nível interno.
\item Usar o teste $t$ ingênuo sobre folds ou sobre sementes.
\item Aplicar Wilcoxon a folds, sementes ou datasets de um mesmo cluster.
\item Concatenar painéis de controles num único MCC.
\item Comparar MCC entre painéis como medida de qualidade dos controles.
\item Afirmar equivalência sem margem predeclarada.
\item Rodar um $K$-fold externo sobre o dataset inteiro, o que reciclaria o teste externo selado.
\end{itemize}

%=====================================================================
\section{Limitações deste relatório}
%=====================================================================

\begin{itemize}[leftmargin=1.4em]
\item As fórmulas de \citet{dietterich1998} e \citet{alpaydin1999} foram verificadas nas formulações publicadas e conhecidas. Apenas a de \citet{bouckaert2004} foi conferida no texto completo durante esta revisão.
\item O comportamento dos testes com MCC e com dados agrupados não foi simulado aqui. Uma simulação com dados sintéticos agrupados, sem resultados reais, seria o passo natural para escolher entre o $t$ corrigido e o $t$ bayesiano correlacionado na análise S1.
\item A margem $\delta$ e os limiares de diagnóstico são exemplos. Fixá-los é uma decisão do autor.
\item Os fatos do projeto refletem os arquivos do repositório em 4 de novembro de 2026.
\end{itemize}

\bibliographystyle{plainnat}
\bibliography{references}

\end{document}
